\chapter{可复现性与材料整理说明}
% \chapter{Reproducibility and Material Organization}

\section{样本与输入版本记录}
% \section{Record of Samples and Input Versions}

% [PENDING]
本附录用于集中记录论文中图表、表格和结果所依赖的输入版本信息，包括数据时期、触发版本、
ntuple 或 skim 版本、MC 配置摘要以及外部分析工作区中的对应标签。

\section{图表来源说明}
% \section{Figure and Table Provenance}

% [PENDING]
由于论文所用图通常在其他工作区生成，本仓库不要求保存作图脚本。
但每张进入论文的图都应在 \texttt{notes/provenance.md} 中至少记录：

\begin{itemize}
  \item 图文件名；
  % \item the figure file name;
  \item 外部来源位置或工作区说明；
  % \item the external source location or workspace description;
  \item 生成日期；
  % \item the production date;
  \item 关键输入版本；
  % \item the key input versions;
  \item 与正文对应的章节位置。
  % \item the corresponding chapter location in the main text.
\end{itemize}

\section{结果填充清单}
% \section{Checklist for Filling Results}

% [PENDING]
最终定稿前，应逐项核对以下信息是否已经补齐：

\begin{itemize}
  \item 数据样本和积分亮度；
  % \item the data sample and integrated luminosity;
  \item 触发与对象选择；
  % \item trigger and object selections;
  \item 效率修正和系统误差；
  % \item efficiency corrections and systematic uncertainties;
  \item 拟合结果或上限结果；
  % \item fit results or upper-limit results;
  \item 正文、图表和参考文献之间的交叉引用。
  % \item cross-references among the main text, figures and tables, and bibliography.
\end{itemize}
